%============================= abstract.tex================================
\chapter*{Abstract}%
\addcontentsline{toc}{chapter}{Abstract}%
\vspace{-25pt}

\begin{spacing}{1.2}
\noindent Access to Indian legal information remains hindered by complex statutory language, fragmented document repositories, and language barriers for non-English speakers. Citizens seeking redress through civic complaint systems face difficulties understanding relevant laws such as the Bharatiya Nyaya Sanhita (BNS), Bharatiya Nagarik Suraksha Sanhita (BNSS), Right to Information (RTI) Act, and Consumer Protection Act. Traditional search interfaces return raw legal text without context, requiring professional legal interpretation. To address these challenges, this project presents \textbf{CiviQ}: an AI-Powered Legal Information Retrieval system using \textbf{Retrieval-Augmented Generation (RAG)}. The system ingests PDF legal documents, extracts text using \textbf{PyMuPDF}, and applies overlapping word-chunking (900 words per chunk, 150-word overlap, 40-word minimum). Text chunks are embedded using \textbf{sentence-transformers (all-MiniLM-L6-v2)} generating 384-dimensional cosine-normalized vectors, indexed in \textbf{Pinecone} vector database with batch upserts.

\vspace{8pt}

\noindent User queries are embedded identically, and top-k semantic retrieval (default k=5) fetches relevant passages. \textbf{Google Gemini 2.5 Flash} generates grounded answers citing retrieved sources. The system detects Malayalam queries via Unicode block analysis (U+0D00--U+0D7F) and translates them to English before retrieval, with response translation back to Malayalam. The production architecture comprises a \textbf{FastAPI} backend (Python 3.12) with endpoints for question answering, document upload analysis, and health checks, serving a \textbf{React 19 + Vite 8 + Tailwind CSS 4} frontend with voice input (browser speech-to-text), PDF upload, language selection (auto/English/Malayalam), top-k slider, source citation display, and document summaries. The corpus includes BNS 2023, BNSS 2023, BSA 2023, Consumer Protection Act 2019, Indian Contract Act 1872, RTI Act 2005, and Constitution of India. The system provides accessible legal guidance for civic complaint resolution and information rights.
\end{spacing}

\thispagestyle{plain}
%=======================================================================
